\documentclass[a4paper, 12pt]{report}




% Some highly suggested packages, please read their manuals.


\usepackage{natbib}
\bibliographystyle{apalike}
\usepackage{url}

\usepackage[utf8]{inputenc}
\usepackage{color}
\usepackage[top=2.5cm, bottom=1.5cm, left=2.5cm, right=1cm]{geometry}
\usepackage[pdftex]{graphicx}
\usepackage{tikz}
\usepackage{breqn}

\usepackage{cleveref}
% Logos
\newcommand{\ulb}{\includegraphics[scale=1.1]{logo_ULB2.pdf}}
\newcommand{\polytech}{\includegraphics[scale=0.35]{logo_polytech_FR.pdf}}

% Polices
\definecolor{ULBblue}{rgb}{0,0.2196,0.5765}
\newcommand{\fontTitle}{\sffamily \Huge\selectfont \color{ULBblue}}
\newcommand{\fontSubtitle}{\sffamily \LARGE \selectfont \color{ULBblue}}
\newcommand{\fontText}{\sffamily \selectfont}
\newcommand{\fontColor}{\sffamily \selectfont \color{ULBblue}}

% Titre
\newcommand{\titleA}{\fontTitle{Première ligne de titre du mémoire}} % Titre identique au titre remis au secrétariat
\newcommand{\titleB}{\fontTitle{Deuxième ligne de titre du mémoire}} % (dans la langue de rédaction a priori)
% Sous-titre
\newcommand{\subtitle}{\fontSubtitle{Ligne du sous-titre du mémoire}}
% Titre du diplôme
\newcommand{\diplomaA}{\fontText{Mémoire présenté en vue de l’obtention du diplôme}} % A laisser en Français
\newcommand{\diplomaB}{\fontText{d'Ingénieur Civil [...] à finalité [...]}}

% Etudiant
\newcommand{\student}{\textbf{\sffamily \large [Prénom Nom]}}

% Supervision
\newcommand{\promAa}{\fontColor{ primary eur}}
\newcommand{\promAb}{\fontText{Professeur [Prénom Nom]}}
\newcommand{\promBa}{\fontColor{Co-Promoteur}}
\newcommand{\promBb}{\fontText{Professeur [Prénom Nom]}}
\newcommand{\promCa}{\fontColor{Superviseur}}
\newcommand{\promCb}{\fontText{[Prénom Nom]}}
\newcommand{\deptA}{\fontColor{Service}}
\newcommand{\deptB}{\fontText{[Nom du service]}}

% Année académique
\newcommand{\yearA}{\fontColor{Année académique}}
\newcommand{\yearB}{\fontText{20xx - 20xx}}

\begin{document}

	\thispagestyle{empty}
	\newgeometry{top=2.5cm, bottom=1.5cm, left=2.5cm, right=1cm}
	\setlength{\unitlength}{1mm}
	\noindent\begin{picture}(175,257)
	
		\put(0,245){\polytech}
		\put(153,139.5){\ulb}
		
		\put(8,155){\makebox(150,10)[l]{\titleA}}
		\put(8,145){\makebox(150,10)[l]{\titleB}}
		\put(8,135){\makebox(150,10)[l]{\subtitle}}
		
		\put(0,75){
		\begin{tikzpicture}[scale=0.1]
		\fill [fill=ULBblue](0,0) rectangle (0.8,90);
		\fill [fill=ULBblue](0,57) rectangle (152,57.8);
		\end{tikzpicture}}
		
		\put(8,120){\makebox(150,5)[l]{\diplomaA}}
		\put(8,115){\makebox(150,5)[l]{\diplomaB}}
		
		\put(8,75){\makebox(150,10)[l]{\selectfont \student}}
		
		\put(8,44){\makebox(80,5)[l]{\promAa}}
		\put(8,39){\makebox(80,5)[l]{\promAb}}
		\put(8,31){\makebox(80,5)[l]{\promBa}} % Commenter la ligne si pas nécessaire
		\put(8,26){\makebox(80,5)[l]{\promBb}} % Commenter la ligne si pas nécessaire
		\put(8,18){\makebox(80,5)[l]{\promCa}} % Commenter la ligne si pas nécessaire
		\put(8,13){\makebox(80,5)[l]{\promCb}} % Commenter la ligne si pas nécessaire
		\put(8,5){\makebox(80,5)[l]{\deptA}}
		\put(8,0){\makebox(80,5)[l]{\deptB}}
		
		\put(145,5){\makebox(30,5)[r]{\yearA}}
		\put(145,0){\makebox(30,5)[r]{\yearB}}
	
	\end{picture}
	\restoregeometry



\chapter{Introduction}
Social network analysis is an interdisciplinary discipline originally developed under the influence of sociology and mathematics \citep{history_social}. It consist of using graph and network theory to represent links among individuals as a network in order to analyze them. \\

Research in computer science are developing semantically-oriented techniques to analyze fiction. \cite{character_country} had presented a method to extract social network from literary which allows to apply Social Network Analysis techniques on it. \cite{movie} presented a method to extract a social network from movies. Automated method that analyze text of fiction also produce metadata that help to analyze them. For instance \cite{original} produce the social network novels but also count the mean distance between two dialogs in the novel. This measure produce meaningful result when it is used to classify novels. Social network can be seen as a metadata with the particularity of being present in any fiction no matter the support of this fiction. Other metadata are only measurable in movies, novels... \\

This master thesis consists firstly in the development of a software that extracts social networks and other metadata from novels and movie scripts. Secondly it consists in the analysis of the topology of the extracted networks. The software has been originally developed during an other master thesis \citep{original_thesis} and further advancement have been achieved in \cite{original}. It was only working on novels. For this reason I will focus on the modification that I have made and the evolution of the state of the art since this period. The global way the software works and the main steps will also be explained more briefly. \\







\chapter{Method of the software}

In this part of this Master Thesis, I will briefly explain what are the main step executed by the software to extract social network and meta data from novels or script. The modification of the program that I've made will be more detailed. Every information about the initial software are available at \cite{original}. The initial software was only analyzing novels, some adjustments have been made to extend it to movie scripts.

\section{Assumption and conception of the social network}
There is multiple way of constructing a social network from a text.  Especially because there is multiple way of analyzing the structure of the narrative. \cite{original} choose to see a narrative as a succession of event. The text is composed of a series of \textbf{conversation} that involve a set of characters. Each conversation can be subdivided into a succession of mini-event that are called dialog. In a \textbf{dialog} a person chosen in the set of character involved in the conversation will address a message to a the other characters from the set, the \textbf{audience}. The person that send the message is called the \textbf{speaker}. Also all sentences that are out of dialogs are called \textbf{context} and linked with the following each conversation or the one that is currently happening. Sentences of context can be located between dialogs from a same conversation. Characters that appear together in a conversation are linked in the social network with a weight that correspond to the number of conversation they are both in.

The narrative is also divided in \textbf{scene} that could contain multiple dialog. The scene represent additional cut where all the previous context should be thrown up. In novels scene are used to represent new chapter, in movie script scene are directly annotated.

\section{Text analyzed}
Most of the novels analyzed in this work have been provided by \cite{original}. They has been selected to represent a variety of popular series and some classic novels. They choose to select series to observe if social network from novels being part of a serie are sharing common properties.  I've added to it the novel ``Pride and Prejudice'' to compare my result of character extraction with the results of \cite{character_meta}.
Various script of movies have been added to compare social network developed in literary fiction and movies. There is both original scenario and adaption from books in order to measure difference between those type of movies. Adaptation of analyzed book have been added when they were available.



\begin{table}
\center
\scalebox{1}{%
\begin{tabular}{|c|c|}
\hline
Novel & Script \\
\hline
\multicolumn{2}{|c|}{Pride and Prejudice}\\
Harry Potter 1$\rightarrow$7 & Harry Potter 1-2-3-4-6-7 \\
A song of Ice and Fire 1$\rightarrow$5&Black Panther\\
The Lunar Chronicles 1$\rightarrow$3 &Thor Ragnarok\\
His Dark Materials 1$\rightarrow$ 3& Blackkklansman\\
The Mortal Instruments 1$\rightarrow$6 & Boyhood\\
The Liveship Traders 1$\rightarrow$3 &Halloween\\
The Wheel of Time 1$\rightarrow$14 & Joker\\
Rain
Wild Chronicles 1$\rightarrow$4 & Jurassic Parc  1$\rightarrow$3\\
&Lord of the ring  1$\rightarrow$3\\
& Shrek 1 and 3\\
&The devil wears prada\\
&Titanic\\
& Alien 1$\rightarrow$3\\
\hline
\end{tabular}}
\caption{}\label{Novels and script whose social network has been extracted. `Pride and prejudice' has been used under the form of the novel and under the form of the adaptation's script. In total there is 46 novels and 25 scripts.}
\end{table}

\section{Implementation}

Most of the NLP function used on this software come from the library pattern \citep{pattern} only available on python 2. As new tools have been developed in python-3 library, the choice have been made to pass the code to python 3. 
In order to keep the structure of the code developed with the tool pattern, a python 3 branch of the library have been used. The branch is still in development and every function are not runnable but I manually fixed the needed function use it on the software.  My fork of the library is available at \textit{https://github.com/antheymans/pattern}.\\

To use new function that doesn't exist on the pattern library, the choice has been made to use the Spacy library \citep{spacy}. This library develloped by the MIT has the advantage of being well-documented. This is completed by the use of NLTK \citep{NLTK} which posses a wide variety of data and corpus for NLP.  A lot of state of the art article use the standford CoreNLP package \citep{CoreNLP} wich also haves all the state of the art tool. However this package is written in java. There exist a python adaptation of it but it is is poorly documented and doesn't include all the tool of the java version \citep{CoreNLPPython}.\\



\section{Formatting of input}
In order to be used as input of the program, a narrative should be on the form of a text and  transformed in order to follow a set of rules. The program have been made to exploit information available in books and movie scripts but a narrative following any other format can be used in the program if it follows the rules.

\begin{itemize}
\item The text should be stored in a \textit{.txt} file with an \textit{utf-8} encoding.
\item All information located in the document should be part of the narrative. The title can be kept but any preface or thanks from the author will be considered as being part of the narrative.
\item All dialogs should be delimited by double-quote. Double quote can not be used for other purpose. All double quote present in the text should be replaced with simple quote.
\item If consecutive dialogs belong to the same conversation and have the same speaker, they should be concatenated. There is an exception to this rule is there is a context between those dialogs.
\item A new scene should be represented by an empty line. No other empty lines are authorized.
\item If the speaker of a dialog is annotated, it should be indicated in uppercase on the previous line. The name of the speaker should be preceded by ``\_'' and followed by  ``:''.\\
 For instance: \\
	\_TINTIN:\\
	``Milou, go outside!''\\
\end{itemize}

The formatting is made manually using regex. A script have been created to transform automatically scripts having the most common structure  but it often requires minor manual adjustment.

\section{Preprocessing of the text}
Initially all the sentences of the text are extracted, parsed and divided into chunks. Annotated speaker are extracted and linked with the related dialogs. Sentence belonging to a same dialog are grouped. The location of scene breaks is also stored in order to later cut the conversations at this points. The sentiment associated to each sentence is also computed using 2 scores: a measure of polarity that say how much the sentence is negative or positive and a measure of subjectivity that say how much the sentence tends to give personal opinion.

\section{Construction of conversations}
To group dialog in conversation, the program should be able to decide if 2 consecutive dalogs from a scene belong to a same conversation. The difficult part is that there is often sentences of context between dialogs of a conversation. The method develloped by \cite{original_thesis} use the size of the spacing between dialogs to perform this task. It build conversations as a a list of dialogs having an associated list of non-dialog sentences. The second list is called the context of the conversation.\\
The program begins by counting the number of spacings separating each dialogs. The spacing between dialogs belonging to different scene is not taken into account. From the number of occurrence of each value of spacing, a threshold value will be chosen. Dialogs separated by a spacing above this threshold value are not part of the same conversation. The threshold value is chosen using the following formula:

\begin{equation}
\begin{split}	
treshold = max\{spacing | frequency(spacing) > 100 \\
\vee (frequency(spacing) >= 10 \\
\wedge frequency(spacing) >= frequency(spacing + 1) \cdot 2 \\
\wedge frequency(j) > frequency(spacing) \forall j \in [1,spacing])\}
	 \end{split}
\end{equation}
The formula has been empirically chosen by \cite{original_thesis}.

A conversation end with the last dialog of the scene or with a dialog being separated by its successor by a spacing bigger than the threshold value. All sentences of context located in a conversation or between the current conversation and the previous one are grouped. This group is the context of the conversation.\\



\section{Character identification}
Character extraction is a major task of social network extraction. ``Character identification consists in detecting which characters appear in the considered narrative, and when exactly they appear in this narrative'' \citep{fiction}. 2 sub-task can be considered as separated from each other:
\begin{enumerate}
\item The extraction of a list of characters.
\item The detection in the text of character appearance.
\end{enumerate}

The first sub-task will be deeper analyzed on this work, as a new detection system have been developed to improve the result from \cite{original} on novels. On this part, scripts and novels need different processing as scripts have annotated speaker for most of their dialog.  


\subsection{Character extraction in Novels}
According to \cite{fiction},  character in novels may appear on the form of a proper noun, a pronoun or an anaphoric noun phrase. Proper nouns that are composed of a single word are called proper names. In this work, character are only detected when they are on the form of a proper noun. This choice is motivated by the fact that characters will be later connected when they are mentioned in the same conversation. The cost of this simplification is the lost of smaller characters that appears only under the form of an anaphora. I considered that in most cases, a character taking part in a conversation is mentioned at least once under the form of a proper noun. In our situation, the first step of \textit{character identification} on a written support is the extraction of names that represents the characters and the linking of aliases (names that refers to the same character). Once all extracted names are bound with a character, each occurrence of a name in the story signal the appearance of the associated character. As the task of extracting a set of proper nouns and binding aliases is error-prone, some authors decided to do it manually \citep{agarwal-etal-2013-automatic}.  \cite{he-etal-2013-identification} proposes to build automatically a list of character from wikipedia but this method has the disadvantage to focus only on main characters and to be dependent on external information unavailable for some stories.  Here we will focus on automatic method because they can be used on many different text with a minimal pre-processing.\\

\subsubsection{The disadvantage of proper names over proper nouns}
A major modification of the program concern this first step of the character identification process. The original process was only considering that a character could be represented by a proper name and was linking names that appears together more than 1 times over 3. But using only proper names over all proper nouns is a simplification that makes the program loose a lot of information. The use of proper nouns makes the binding of nouns more complex but also allow to use more information's to perform this binding. The current state of the art contains methods to perform this task. \\

 \begin{tabular}{lll}
\textbf{chunk} & \textbf{names extracted using proper nouns} &  \textbf{names extracted using proper names} \\
   Ron and Hermionne & Ron AND Hermionne & Ron\\
   you Mr. H. Potter & Mr. H. Potter &  Potter\\
   dear Harry Potter & Harry Potter & Harry OR Potter \\
   James Potter& James Potter & James OR Potter \\
   yer brother Charlie & Charlie & Charlie\\
\end{tabular}


 The fact that proper names are only composed of a single token makes a perfect binding of names designating the same character impossible. The pairing is needed to link first names and last names of characters but in some cases multiple proper nouns that should not be linked appear in the same chunk. Here are some common mistakes in character extraction with example from the method of \cite{original} used with the book "Harry Potter 1".
\begin{enumerate}
    \item \textbf{When both the first name and the last name of a character are used to identify a character, they are not paired}: "Harry" and "Potter" are labeled as different character.
    
    \item \textbf{When multiple characters share the same first name or last name, if the linking is made they will all be clustered as a single character}: the names referring the 8 characters of the "Weasley" family are mixed in 4 cluster. Each cluster is composed of aliases referring to 2, 3 or 4 characters.
    
    \item \textbf{Some characters are referred with proper nouns that should not be used for pairing}: "Mr" is paired with "Dursley" and "Weasley" because the text often refer to "Mr Dursley" or "Mr Weasley". It causes the clustering of unrelated words. The same problem appear with other title like "professor". 
    
    \item \textbf{Some chunks contain multiple characters}: The algorithm link to a single character the chunks "Ron and Hermionne", "Mr and Mrs Weasley" or "Fred and George".
\end{enumerate}

\subsubsection{Name Entity Recognition}
The task of labeling group of words as Entity is called Name Entity Recognition (NER). Those entity includes \textit{person}, \textit{location} and \textit{organisation} \citep{libraries}. This is the most common way to extract characters names from a novel.
The best NER methods are supervised or semi-supervised and trained with annotated datasets of news, text from social media or biomedical data \citep{NER_survey_recent}. This decreases their performance on novels, especially novels from fantasy or the older ones \citep{NER}.  Unsupervised methods typically rely on rules and domain-based knowledge, it makes them completely domain dependent.  POS-tagger may be considered as naive program of NER. The original software \citep{original} was extracting all words labeled as proper names by a POS-tagger with the major issue of considering only single-word names while in \cite{quoted}, proper names are tagged using a POS-tagger and contiguous proper names are considered as a proper noun. \cite{character_meta} also extracts subjects of verbs present in a dataset of verbs strongly associated with ``person'' entity. This techniques allows to detect anaphoric nouns and not only proper nouns. \cite{ structure_clustering} store the number of times that words are classified at location or person, so words that are likely to be location are removed from the location list. \\
Pattern the NLP  library of python that is used in the software doesn't have any NER module. Spacy, another NLP-library of python have an available module that performs NER using `Conditionnal Random Fields', a statistical model that uses supervised learning. Pre-trained model are available but the model could also be trained manually with an annotated dataset.  The most common tool for NER is the Stanford Named Entity Recognizer (NER) wich is only available on java. For reason of practicability, it has not been considered in this work.


\subsubsection{Unification of Character Occurrences}
Unification of Character Occurrences is the task of unifying all mentions of a same character in a narrative \citep{fiction}. When only proper nouns are considered, this task can be simplified into Alias Resolution: the linking of all names that refers to a same characters and the making a differentiation between names that refers to different characters \citep{book_social}.\\
The binding of names can be done using a measure of string similarity, set of rules or using meta-information of strings such as an inferred gender. In multi-stage methods, proper nouns are are divided between cluster, each of them being associated to a single character. The clusters are merged following a sequence of conditions. Multiple methods have been developed to solve this task with their own specificity without that one method stand out from others.\\
The original software \citep{original} was linking proper names that appears together 1/3 of the times to link first name with last name with the major inconvenience of binding different characters sharing a same first name or last name and leaving unbound name and lastname of characters that are more frequently called by only one of those words. 
However most of the state of the art methods focuss on proper nouns: \cite{delete5} discards all proper nouns that appears less than 5 times, then try to bind multi-words nouns between them before binding them with single words nouns. \cite{structure_clustering} also classify names into multiple classes and apply different rules on those classes. \cite{character_meta} propose to bind nouns that share words except in some cases, such as nouns sharing a last name but having a different first name. \cite{variation} proposed a method to bind entity representing art object by generating variations of proper nouns and linking them with names corresponding to those variations, this method have been applied on characters detection in novels by \cite{quoted, character_meta}.  In \cite{delete5, structure_clustering, quoted},  a gender is inferred from each proper nouns using a list of masculine and feminine first names and gendered titles to avoid to merge cluster of names having different genders.\cite{ structure_clustering} a sliding window is also used on the text to detect pronoun near nouns and use it to infer the noun gender. It ends by removing infrequent characters, characters whose the total number of apparition of the cluster of names  is smaller that a threshold. This should reduce the number of ``false positive'' characters (cluster of names that are not related with a character), at the cost of minor character. \cite{delete5} draws the conclusion that all methods are error-prone and that the difficulty to obtain annotated data on this task makes any comparison between the different methods very difficult. Even with a dataset containing all characters present in a narrative, the presence of multiple aliases would makes the evaluation of the character extraction very difficult.  \\

Some methods also use co-reference resolution tools to link characters mentions between each others \citep{character_meta}. Coreference resolution tools are program that automatically clusters mention in text that refers to the same entity. They do it using neural networks and are especially used to link names with pronouns or anaphoric noun phrase \citep{coref_deep, coref_deep2} . This is useless to the software here as those form of references are not used in this work. \\

\subsection{Proposed method of character extraction in novels}
As explained in the previous section, the method here extract proper nouns, the extracted nouns are composed of multiple words including honorifics. Also only proper nouns are extracted, anaphoric noun phrase are ignored. The method consist of a multi-pass algorithm that passes 2 times over chunks to extract a list of nouns,  classify them and then pass 5 times over the set of nouns to bind aliases.
\subsubsection{Proposed system of Names Entity Recognition}
\label{section:NER}
To find a NER method for characters extraction in novels, we relied on some hypothesis. Names from fiction have the particularity to be invented in order to be easily recognizable by the reader. We can rely on the hypothesis that in most novels, there is no namesake person. An exception to this hypothesis is the characters ``Barty Croupton'' in Harry Potter, but the younger is called ``Barty Croupton Junior'' to avoid any confusion from the reader. Also we consider that characters from novels should be labeled as person, so we are only looking at the  NER system capacity of distinguish person from other entities. We will later accept some entities as ``speaking object'' as characters. But it has been considered that most of them have specific nouns that are likely to make a NER system to label them as character and the remaining ones should be minor in the narrative. So if we miss one of them, it's not considered as a major issues\\

As explained in the state of the art, a  python library, Spacy, provide a tool for the task of name entity recognition. However this tool has not been trained specifically with novels nor for the task of extracting characters names from a text. It doesn't detect honorifics and works as a ``black box'' which gives sometimes incorrect result. For instance in the book ``Harry Potter 1'' the tool detect has characters: 'Harry bellowed': 1, 'Harry anxiously': 1, 'knew yeh didn',' baker', `yer meddlin','Happy Birthday Harry'. It seems that interjection such as ``yer'' or ``yeh'' are interpreted as nouns and  a lot of verbs or adjectives are seen as part of nouns like with ``Harry bellowed''. Also the tool doesn't make a difference between proper nouns and anaphoric noun such as with ``baker''. This is an issue as anaphoric noun phrase loose their meaning out of their context and considering them as character names could lead to major confusion. For instance the anaphoric noun phase ``the father'' could represent a lot of characters according to the context. For al those reasons it has been chosen to avoid to work with this library. \\

A Name Entity Recognition method has been designed to extract characters names from novels. This technique uses a POS-Tagger to extract proper names from the text and several collection of proper names commonly used in English to isolate characters names. Then proper nouns are extracted from chunks containing proper names. The collection of words contain a list of countries, nationalities, honorifics (classified following the related gender), stop-words, profanities, words related to time and words related to the academic domain. In this method, a  proper name is considered as \textit{valid} if it begins with a capital letter but is not composed with capital letters only. The assumption has been made that proper names contains more than one letter.\\
First of all the POS-tagger is passed on the head word of all chunks to extract  proper names. The POS-tagger used is the one developed by the python library pattern which classify  proper names whether they are related to a location or not. If a word is not labeled as a location, the algorithm considers the word as a proper name if it doesn't appear in the lists of English  proper names and is not the first word of a sentence. First words of sentences are not kept as they are capitalized and so they are more likely to be wrongly classified as proper names. The assumption is made that proper names will appear several time and that at least one of those appearance will not be at the begin of a sentence. Proper names and  proper names related to locations are saved with the number of their appearance.\\
Secondly the algorithm check that words labeled as location don't appear in the list of proper names. In such cases, the word will be removed from the list of proper names if it appeared more often as a location than as a name.\\
The final step of the Names Entity Recognition method consists of extracting proper nouns from proper names. All chunks containing a proper name are isolated.   Then all sequence of consecutive valid  proper names are extracted and considered as proper nouns. For instance from the chunk ``Mr Harry Potter and Ron'' 2 proper names are extracted: ``Mr Harry Potter'' and ``Ron''. The proper nouns are registered with the number of their appearance.\\

In this method, many restrictive conditions are putted on proper names extraction. But the way names are extracted implies that if a name is extracted one time, he will be considered each time as a proper name. For this reason the addition of new names should be very cautious. The hypothesis is made here that a frequent character's name will appear at least once in a situation that allows to univoquely identify him as a character's name. On the original program, only few conditions are observed and result shows that most of the extracted characters don't refer to actual character of the novel. This observation will be explained more formally in the section \ref{performance}

\subsubsection{Proper nouns classification}
\label{section:classification}
All the extracted proper nouns are classified following their gender and their form to make easier the alias resolution. The classification of genders is also used to produce gender-related information about the topology of the network.\\
Firstly all the names are parsed using the python tool \textit{HumanName} from the library {NameParser}. It separates words contained in a noun in a list of honorifics, first name and last name.\\

To infer the gender, the algorithm will used the list of honorifics classified by gender and a list of gender related first name. The gender of nouns containing a first name from the list of first name or an honorific from the list of  labeled honorifics will be assigned. If multiple words in the noun are related to different genders, the majority vote is used. After this phase all nouns are labeled as \textit{masculine}, \textit{feminine} or {neutral}.

Then a category will be assigned to each words following their forms:
\begin{enumerate}
	\item The first category contains nouns that have at least an honorific, a first name and a last name.
	\item The second one contains nouns with a first name and a last name but no honorific.
	\item The third category contains nouns with an honorific and a first name but no last name.
	\item This category contains nouns with an honorific and a last name but no first name.
	\item The fifth category contains the remaining nouns.
	 \item If the first name of a noun is not in the list of proper noun and doesn't respect the criteria to be labeled as proper noun, the noun is considered as an error and is removed.
\end{enumerate}

\subsubsection{Alias resolution}
\label{section:alias}
The alias resolution process consist of a multi-pass algorithm that firstly creates  primary link and secondary link between all  nouns, then  transforms secondary link into  primary link in some cases. After this linking phase, cluster are  build to represent character. Each cluster consists of nouns connected between them by a path of primary link. The most used noun of the cluster is called the head of the cluster and is used to represent the associated character in the social network. In the following explanations, I will speak about ``compatible genders'' between 2 nouns, it means that the gender assigned to those nouns doesn't prevent them to point the same character. Concretely they have the same gender or one of them is neutral. To be connected, nouns have to have compatible genders. Once two nouns are linked by a  primary link, if one of them was gender neutral, its gender will be inferred from its neighbor. In such situations the neutral word will delete all its connections with words of the opposite gender. In the following explanation I will also speack about "primary-connection" or secondary"neighbors"to speak about names connected using a path of primary link and names neighbors from each other by a secondary link. \\

During the first pass, all first name used by nouns of class <3  are taken one by one. Diminutive and nicknames associated to a first name are looked in the set of proper names. Then all proper nouns that possess that first name, a diminutive or a nickname are grouped. In some cases the first name, the diminutive or the nickname has been labeled as last name in other nouns, those nouns are also added. For each pair of nouns from this group that have compatible genders, the nouns are linked
\begin{itemize}
\item
	 using a  primary link if:
	\begin{itemize}
	\item they have the same last name and their first name is the original first name, a diminutive or a related nickname. This also works with nouns that don't have last name as they last name is considered to be the empty string.  For instance: ``D Dursley'' and ``Dudley Dursley''. 
	\item if both of the names have an empty last name or first name. In this case the algorithm consider that they share a same first name/last name even if it is possibly labeled differently in each noun. For instance:``Mr Flitwick'' and ``Professor Flitwick''.
	\end{itemize}
	Once 2 nouns are connected with a  primary link, if one of them was neutral, he will receive the gender of the other word.
\item
	 using an  in primary link if:
	\item  their first name is the original first name, a diminutive or a related nickname and only one of them have an empty last name. 
	\item if one of them have an empty first name.
	Secondary  link represent a potential connection between names which could be unrelated. For instance ''Mr Durlsey'' could be related to ``Dudley Dursley'' or to ``Vernon Dursley'', so he will be linked using a secondary link to both of them. ``Harry Potter'' and ``James Potter'' should not be linked as it appear clearly that this is 2 different character sharing a last name.
\end{itemize}

In a second phase, all last name are taken one by one. Last name that has also been used in previous pass as a first name are not considered because they are likely to be incorrectly labeled. All nouns sharing the last name or a diminutive are grouped. For each pair in this group, if they have compatible genders, they are linked using:
\begin{itemize}
\item a  primary link if both of their firstname are empty. 
\item an in primary link in other cases.
\end{itemize}

Once these phase have been completed, we will observe secondary link and separate those who are due to a true connection between 2 names and the links that should be removed.  We will takes the nouns that contains the less information and try to link them with the nouns that are more likely to be their neighbors. The chosen neighbors is the most common noun among the neighbors. This technique will result in some errors but their numbers should be limited.
\\
The third pass focuses on neutral nouns of category 1 and 2. We will try to find them a gendered alias among their potential partners. Firstly each of them will be grouped with all nouns connected to him using  primary links. Then a second group will be formed with all they secondary-neighbors. Iteratively the most common neighbors  will be removed from the second group and linked to the first group using a  primary  connection. If the noun is gendered, the gender will be assigned to all word from the first group and the process will end. If the noun is also neutral, all nouns  primary ly connected to it will be added to the first group and again the most common noun from the second group will be chosen to be linked with the first group. 
\\
The fourth and fifth pass focus on words from the 3-4-5 categories that are not connected with nouns from the 1-2 categories. Those words consists of a single first name or last name and should be linked to the most common character having this name. Again in the forth pass they are grouped with the primary-connected nouns. The most common noun among their secondary neighbors of 1-2 categories is added to the group and  primary ly linked with it. But nouns that are not linked to any 1-2 nouns will remain unbound. In the fifth pass, all secondary links between those names will be iteratively transformed into primary link. Link between nouns that becomes incompatible due to their gender during the process, are removed.
\\
In the last part of the algorithm, cluster are made with words primary-connected and the most common noun of each cluster is chosen as  head of the cluster.

\subsubsection{Performance}
\label{performance}
As explain in the state of the art, there is no straightforward way to measure the performance of the character extraction techniques. Firstly the performance of the NER and the alias resolution system are not independent: For instance a NER method that add redundant adjectives to nouns will only decrease the performance of the overall social network if the alias resolution system is not able to link all those nouns. Secondly most method from the state of the art don't have open-access implementations. Thirdly there exist almost no annotated data that would allow to automatically score different method. The goal of the algorithm developed here is to find different aliases referring to each character. But there exist multiple way of doing it that are not strictly better from each others. The name ``Weasley'' for instance could be associated accurately to multiple characters in ``Harry Potter'' while associating it with the noun ``Malefoy'' is clearly a mistake.

For those reasons, the method developed in this master thesis will only be compared with the previous method used in the software \citep{original} and with a method of the state of the art, \cite{character_meta},  whose authors provide the result of their character extraction for the novel ``Pride and Prejudice''. I will manually count the number of clusters headword that refer to an actual character to count the number of false positive of character detection. I will also compare the set of main characters detected by our program and the set of main characters detected by \cite{character_meta} and manually check which errors appears. This comparison is not possible with the original software as the program didn't classified the cluster of characters names by number of appearance.\\


To count the number of  false character, I've decided to consider only the most common name of a list of aliases. So a character is considered as false is the most common name is false even if other aliases are valid to designate a character. The name should be part of the canonical name of a character or designate it unambiguously. In the novel ``Harry Potter'' appears a lot of magic object or animals. Animals and speaking object are considered as character as they are potentially part of a conversation. For instance the word ``Hat'' (capitalized) is considered as designating univoquely  a character as it can only refers to a speaking object ``the Magic Hat'' but the words ``Mommy'' or ``Chaser'' could refer to a lot of characters according to context, they are considered as an error. Word related to name, like ``Potter'' are still considered as valid even if they are not univoque. This measure don't give indication of how well alias are bound bind the algorithm but it allow to measure if the algorithm is able to isolate names referring to character in the text. 

\begin{table}
\center
\caption{Number of detected character that refers to actual characters(true positive).}\label{tab1}
\scalebox{1}{%
\begin{tabular}{c|c|c|c|}
\hline
book &Porposed method & \cite{original} & \cite{character_meta}\\
\hline
harry potter 1& 65\% (70/108) & 32\% (61/193) & // \\
The Lunar Cycle & 67\% (24/36)& 26\% (25/96)& // \\
Pride and prejudice &  89\% (44/56) & 46\% (53/116)  & 85\%(61/72 )  \\
\hline
\end{tabular}}

\end{table}


The result shows that the proposed algorithm is far better than \cite{original} at the task of recognizing characters entities in a text of
 novel, for the 3 tested novels, the rate of ``true chatacters'' is each time at least two times bigger. We also see that the result variate
 a lot from one novel to another. This could be explained by the fact that some novels as ``Harry Potter'' develop a large universe and invent a lot of proper names to describe it. Those unknown proper words are particularly difficult to dealt with for the extractor. A program using more syntactic information could may be decrease the importance of this issue. The presented algorithm obtain result very near from \cite{character_meta} however the type of errors that appears on the novel `'Pride and Prejudice'' differs. In \cite{character_meta} most of the errors comes from anaphoric noun phrase such as \textit{servant},  \textit{housekeeper},  \textit{owner} or \textit{assistant}. In our program those element should not be extracted as they are not univoquely linked to a character and depend of their context but we could imagine that an other program that constructs a social network in a context-dependent way could use those information. Most of the errors that appears in our own  program are nouns of location wrongly considered as character: \textit{Netherfield Park}, \textit{Pemberley House}, \textit{Brighton}, \textit{North}. This different of topologies between errors of the 2 considered method could be due to the fact that we develop our own ``Name Entity Recognition'' program adapted to characters recognition that is mainly focus on the form of words wich allows to separate anaphoric noun phrase from proper nouns but can not use syntactic information to separate name of characters from uncommon location. The ``location'' tag of the TAG-Poser is used as a list of countries but it only allows to detect the most common location. On the other hand \cite{structure_clustering} used a common Name Entity Recognition tool that should be able to label nouns as person our location. For future works, it would be interesting to combine rules used by this method with a standard NER tool that would separate characters names from locations. \\

The second test that we will make here will be to compare the main clusters of nouns that represent the characters for the presented method and \cite{character_meta}. This test has the disadvantage of being manual and qualitative more and not automated and quantitative. But not metrics exist to measure the ability for an alias resolution system to connect the related nouns between them.\\

\begin{table}[]
\scalebox{0.65}{%
\begin{tabular}{llll}
False Positive &
  Appearance &
  Gender &
 Aliases \\
0 &
  920 &
  -1 &
  'Elizabeth, Lizzy, Eliza, Miss Elizabeth, Miss Elizabeth Bennet, Bennets, Mrs. Bennet, Mrs. Bennet\\
  & & &'--My Miss Bennets, Miss Bennet, Elizabeth Bennet, Miss Eliza Bennet, All Elizabeth, Miss Eliza, Miss Lizzy \\
0  & 359 & 1  & 'Mr. Darcy, Mr. Darcy:--but, Mr. Darcy--that Mr. Darcy, ", Darcy      \\
0  & 273 & -1 & 'Jane, Miss Jane Bennet                         \\
0  & 212 & 1  & 'Mr. Bingley, Bingley                                  \\
0  & 188 & -1 & 'Lady Catherine, Kitty, Right Honourable Lady Catherine, Catherine \\
0  & 180 & 1  & 'Wickham, Mr. Wickham, Mr. Wickham--when, George Wickham, George    \\
0  & 158 & -1 & 'Lydia, Miss Lydia Bennet, Miss Lydia,                         \\
0  & 148 & 1  & 'Mr. Collins, The Collinses, Collinses,                         \\
0  & 120 & -1 & 'Charlotte, Lady Lucas, Miss Lucas, Miss Lucases, Charlotte Lucas \\
1 &
  114 &
  1 &
  'Netherfield, Netherfield Park, Rosings Park, Park, Rosings, The Netherfield\\
0  & 106 & -1 & 'Miss Bingley, Mrs. Bingley, Caroline Bingley,Bingleys, Caroline \\
0  & 94  & -1 & 'Longbourn, Mrs. Long, Long,The Longbourn                \\
0  & 86  & 1  & 'Mr. Bennet, The Bennets, Bennet                            \\
0  & 57  & -1 & 'Mrs. Gardiner                                        \\
0  & 56  & 1  & 'Sir William, Sir William Lucas, Lucas,Lucases, The Lucases\\
0  & 56  & -1 & 'Mary, Maria, Maria Lucas,                                       \\
1  & 54  & 1  & 'Pemberley, Pemberley House, Longbourn House,                    \\
0  & 54  & -1 & 'Miss Darcy, Mrs. Darcy, Georgiana Darcy,Georgiana              \\
0  & 53  & 0  & 'Meryton, All Meryton                                     \\
0  & 40  & 1  & 'Mr. Gardiner, Gardiner, Gardiners, The Gardiners                \\
\end{tabular}
}
\caption{20 main character detected in ``Pride and Prejudice'' by the presented method. The first column tell if the character refers to a true character from the novel and was used in the previous measurement. The second column contain the total number of appearance of names related to the character and the third one contain a 1 if the infered gender is masculine, a 0 if it is neutral and a -1 if it is feminine. The last column contain all aliases of the character. }

\scalebox{0.7}{
\begin{tabular}{lll}
False positive & Appearance & Aliases                                      \\
0             & 752      & Eliza, Elizabeth, Lizzy, Miss Eliza, Miss Eliza Bennet, Miss Elizabeth, Miss Elizabeth Bennet, Miss Lizzy \\
0             & 313      & Bennet, Mr. Bennet                                                                           \\
0             & 309      & Colonel Fitzwilliam, Fitzwilliam, Mr. Darcy, Mr. Fitzwilliam Darcy                             \\
0             & 291      & Jane, Miss Jane Bennet                                                    \\
0             & 217      & Catherine, Honourable Lady Catherine, Kitty, Lady Catherine, Lady Catherine de Bourgh, Miss de Bourgh\\
0             & 150      & Mr. Collins                                        \\
0             & 115      & Mr. Bingley                                                               \\
0             & 107      & Lady Lucas, Maria, Maria Lucas, Mary, Mary King, Miss King, Miss Lucas                             \\
0             & 73       & Mr. Wickham                                      \\
0             & 58       & Mrs. Gardiner                                                                     \\
0             & 45       & Sir William, Sir William Lucas, William Goulding                                          \\
0             & 39       & Colonel Forster, Forster                                                      \\
0             & 36       & Gardiner, Mr. Gardiner                                                          \\
0             & 34       & Lydia                                                                                   \\
0             & 29       & Mrs. Collins                                                                             \\
0             & 21       & Mrs. Hurst                                                                            \\
0             & 20       & Mrs. Phillips                                                                            \\
1             & 18       & servant                                                                                 \\
0             & 16       & Caroline                                                                                  \\
1             & 16       & housekeeper                                                                    \\
\end{tabular}}
\caption{20 main characters detected in ``Pride and Prejudice'' by \cite{character_meta}. The first column tell if the character refers to a true character from the novel and was used in the previous measurement. The second column contain the total number of appearance of names related to the character. The last column contain all aliases of the character. }

\end{table}

In the character extraction of Pride and Prejudice, a first difficulty that appears is the extraction of the Bennet family. In this family there is 5 daughters named in the text and a mother only referred as Mrs. Bennet. In the presented method, the character of Mrs.Bennet have been mixed with ``Elizabeth Bennet'', the most common character among the sisters. About that last characters, both method have been able to link it with her nicknames ``Eliza'' and ``Lizzy'' . In the Bennet family both method isolate correctly the sisters ``Jane''' and ``Lydia'' but \cite{character_meta} didn't record the alias ``Lydia Bennet''. Our method have a second error, the character `` Mr. Fitzwilliam Darcy `` is separated in 2 characters: ``Mr Darcy'' and ``Mr Fitzwilliam'' that don't appear on the figure as it is less used. But the character ``Charlotte Lucas'' is correctly extracted while \cite{character_meta} mix it with an other character`Mary King'. They also mix the character ``William Goulding'' with the character ``William Lucas''. We see that they have considered 2 anaphoric noun phrase as character: ``servant'' and ``housekeeper'' while our method have wrongly considered 2 location as character ``Permberley House'' and ``Netherfield Park''. Our method bind family names such as ``'The Gardiners'' with the most common character of the family while the other method avoid to use them.\\

The observation of the characters extracted by both method shows some tendencies. Our method tend to extract more aliases but also extract some incorrect ones. It also bind set of aliases more cautiously with result in dividing a character in two set of aliases but avoid to mix between them completely different characters. The only case of mixed characters in our method comes from characters of a same family that differs only by a title, both associated to the same gender. As previously shown, the proposed method ignore anaphoric noun phrase but is likely to consider locations as characters. This analysis don't allow to consider one of the method to be really better than the other. From this we can consider that our method is approximately at the level of the state of the art.\\

\subsubsection{Further improvement}
The method develloped here for character extraction in novels gives satisfying result but could have several improvements. Firstly a NER tool should be used in order to use syntactical information to separate location from characters names. Secondly some rules could be added to link aliases between them for instance in the case of characters having multiple last names. Then a way of using anaphoric noun phrase in the system could be used in order to detect unnamed characters but this addition should be handle carefully as the addition of anaphoric noun phrase could lead to the extraction of many nouns that are not univoquely linked to a character.

\begin{table}
\center
\scalebox{0.7}{%
\begin{tabular}{|c|c|}
\hline
Speaker & Issue\\
\hline
_BEAUREGARD - KLAN NARRATOR (O.S.) & The name is given with an anaphora  \\
_JEROME TURNER (V.O.)(CONT'D)& additional information are given with the annotation \\
_RON (CONT'D) / _RON STALLWORTH& This is 2 different annotation that refers to the same character\\
_RON STALLWORTH, FLIP AND JIMMY& An annotation can refer to multiple character\\
_CSPD OFFICER BRICKHOUSE & Character can be represented by their function and not only by their name\\
_AMERICAN TERRORISTS & Groups are also considered as speaker in some situation\\
\hline
\end{tabular}}
\caption{Speaker annotated in the script of the movie Blackkklansman}\label{annotation}
\end{table}

\subsection{Character extraction in scripts}
As previously explained, in all scripts, the speaker is directly annotated for most of the dialogs. The list of all annotated speaker is a good approximation of the character list. However it causes a list of issue that I will detail here under. Example annotation are given on table \ref{annotation} to illustrate issues caused by using those annotations.\\



\begin{itemize}
\item 
 The observation of scripts shows that those annotation mostly use canonical characters names (with the first and last name for the character). But sometimes character are only referred with their first or last name. So a minimal alias resolution system is needed. 
\item
As the name of characters used for annotation are mostly canonical name, the list will miss aliases of the characters. Especially nicknames and familiar names used by characters that are very intimate with each others.
\item If a character don't talk during a conversation and is only mentioned by other characters using unknown aliases, he will not be detected as taking part of the conversation. The network may miss some relations.
\item 
If a character character is part of the audience in some conversation but never talk, he will not be added in the social network. We may loose some minor character.
\item Information could be difficult to extract from annotations. Multiple characters or extra information could be written and everything is uppercase which makes difficult to extract names.  A system filtering the information is needed.
\item As script are not perfectly formatted, the automatic labelisation of speaker and dialogs using regex may lead to errors. I've manually corrected some of those but I can't look all the text for them. This causes the incorrect extraction of words as characters. 
\end{itemize}

In order to use annotations while diminishing the number of related errors, a method has been developed. The method is inspired from the previous method for novels with many simplifications. The method is develloped here under:
\begin{enumerate}
\item In pre-processing: When a parenthesis is detected, we assume that the end of the sentence is an additional information which is not kept. All words are capitalized. The sentence is spitted in multiple speaker when ``And'', ``/''  or ``,`` are detected.   Useless characters such as ``;'' or ``.'' are removed.
\item At the begin of character extraction, all speakers are registered as proper nouns.
\item Each of the proper names that composed the speakers are also registered as proper nouns and linked with the corresponding speaker using a secondary link.
\item Speaker are also splitted when a ``-'' is detected. ``-'' are used in scripts to gives multiple name to a same speaker, for instance: ``_BEAUREGARD - KLAN NARRATOR ``. All names generated from the split are considered as proper nouns and link with the corresponding speaker using a primary link.
\item The genders of all nouns are infered using the technique explained in section \ref{section:classification}.
\item For all extracted nouns that have been at least one time extracted after a split (they have been part of a bigger speaker in one occurrence), the  neighbor having the most appearance is chosen. If the neighbor have a compatible gender, he is linked with a primary linked and their gender is equalized following the technique of section \ref{section:alias}.
\item Cluster are made of nouns connected by primary links. The most common of the nouns is the head of the cluster. Each cluster represent a character and its aliases.
\end{enumerate}

The method allow to increase the number of detected aliases, link aliases between them and infer gender from names. It also allow to filter useless information from speaker and split multiple speakers when their appears. 	As speakers are not always properly chosen in pre-processing, we may still have some errors. There is also still the problem with non-speaking character. However those characters are minor and we consider that it is not an important issue.  On the movie ``Blackkklansman'', no errors are detected. All the extracted characters are really speakers from the story and their aliases are correctly bound. However some of the speakers are not human being: ``Pre - Recorded Message'' refers to a sentence given by the messagery of a phone, ``All'' refers to a sentence said by a crowd. I remind that in section \ref{performance}, I've defined speaking object as potential character. The list of extracted characters of ``blackklansman'' is given on figure \ref{character_bkm}.\\

An other solution would have been to use the character extractor developed for novels on script. The table \ref{tab1} shows that this system is  doing many errors. Using this system on script could lead to even more errors as script have a less regular format than novels. This difference could be explained by the fact that scripts are not commercial product, they are not salable so error in the format are not corrected while the script is understandable for the users. In particular we found more uppercase words that makes the extractions of proper nouns more difficult. However the use of techniques associated to novels on scripts could improve the result for further advancement when character extractor will decrease this number of false positive, as techniques associated to novels use more information than just the annotation.\\

\begin{table}[]
\scalebox{0.7}{%
\begin{tabular}{|l|l|l|l|l|l|l|}
\hline
Cluster mentions & Gender   & Main Name                  & Other names      &           &          &            \\ \hline
312              & 1        & Ron Stallworth             & Stallworth       & Ron       &          &            \\ \hline
158              & 0        & Flip                       &                  &           &          &            \\ \hline
120              & 1        & Felix                      &                  &           &          &            \\ \hline
78               & 1        & Patrice                    &                  &           &          &            \\ \hline
58               & 0        & Chief Bridges              & Bridges          & Chief     &          &            \\ \hline
57               & 1        & Walter                     & Walter Breachway & Breachway &          &            \\ \hline
57               & 1        & Devin Davis                & Davis            & Devin     &          &            \\ \hline
29               & 1        & Connie                     &                  &           &          &            \\ \hline
29               & 1        & Walker                     &                  &           &          &            \\ \hline
25               & 0        & Ivanhoe                    &                  &           &          &            \\ \hline
23               & 0        & Sgt Trapp                  & Trapp            & Sgt       &          &            \\ \hline
20               & 0        & Landers                    &                  &           &          &            \\ \hline
20               & 0        & Kwame Ture                 & Ture             & Kwame     &          &            \\ \hline
19               & 1        & Beauregard - Klan Narrator & Klan Narrator    & Narrator  & Klan     & Beauregard \\ \hline
19               & 1        & Jimmy                      &                  &           &          &            \\ \hline
13               & 1        & Jerome Turner              & Turner           & Jerome    &          &            \\ \hline
11               & 1        & Mr Turrentine              & Turrentine       & Mr        &          &            \\ \hline
10               & 0        & Agent Y                    & Agent            &           &          &            \\ \hline
5                & 0        & Officer Mulaney            & Mulaney          & Officer   &          &            \\ \hline
5                & 0        & Pleasant Man               & Man              & Pleasant  &          &            \\ \hline
5                & 1        & Hakeem                     &                  &           &          &            \\ \hline
3                & 0        & Wheaton                    &                  &           &          &            \\ \hline
3                & 1        & Butch                      &                  &           &          &            \\ \hline
3                & 1        & Jesse                      &                  &           &          &            \\ \hline
3                & 0        & Cspd Officer Brickhouse    & Brickhouse       & Cspd      &          &            \\ \hline
3                & 0        & Cspd Officer Myers         & Myers            &           &          &            \\ \hline
2                & 0        & Black Mass                 & Mass             & Black     &          &            \\ \hline
1                & 0        & Officer Cincer             & Cincer           &           &          &            \\ \hline
1                & 0        & Pre - Recorded Message     & Recorded Message & Message   & Recorded & Pre        \\ \hline
1                & 1        & Ron Stallworth             &                  &           &          &            \\ \hline
1                & 0        & All                        &                  &           &          &            \\ \hline
1                & 0        & Voice                      &                  &           &          &            \\ \hline
1                & 0        & Klansmen                   &                  &           &          &            \\ \hline
1                & -1       & Odetta                     &                  &           &          &            \\ \hline
1                & 1        & Josh                       &                  &           &          &            \\ \hline
1                & 0        & American Terrorists        & Terrorists       & American  &          &            \\ \hline
\end{tabular}}
\label{character_bkm}
\caption{Extracted characters for the script of the movie Blackkklansman. The first column contain the total number of times each alias of the character have been recorded as speaker. The second one contain the gender of the character. Then comes all aliases of the name. }
\end{table}

\subsection{Identification of character appearance}
\label{identification}
This step consists in binding each occurrence of character mention to the corresponding character and to identify speakers and listeners in each dialog. The algorithm have been mainly designed by \cite{original_thesis} but some adaptations have been made. The previous algorithm was using proper names and chunks associated to this proper names. The new version use proper nouns only.\\

\subsubsection{Detection of potential speaker and listener}
The first phase of this task is the identification of potential speaker and listener in dialog. This don't require the use of the character list previously extracted. It can also be considered as a part of the conversation construction.\\

The program already parsed all sentences and separate dialogs from context. The parsing include the tagging of words by a POS-tagger. All group of words labeled as subject by the POS-tagger are considered as potential speaker. Words labeled as object of the sentence are potential listener. Those group of words are mainly pronouns and anaphoric noun phrase, that are useless for the social network construction. For this reason, proper nouns will be extracted from those group of words. If no proper noun can be extracted, the subject or object is not kept. The extraction follows the step explained section \ref{section:NER}. The initial software was keeping chunks but the system has been modified to be compatible with the new way of extracting characters.

\subsubsection{Filtering of potential speaker and listener}
\label{filtering}
For each conversation, each character detected as speaker or listener need to be associated to a character extracted in the text. I remind that potential speaker and listener are proper nouns. To look for a character, the following step are followed:\\
\begin{enumerate}
\item If the noun correspond to an alias, they are linked.
\item If it is not directly related to an alias, the nouns is separated into its different words. We link the noun with each alias that correspond to one of those words.
\item At the end of the previous steps, the noun is linked with 0, 1 or multiple aliases. The characters corresponding to each aliases will be looked up and we consider that each of them made an appearance at  this sentence.
\end{enumerate}

\subsubsection{Identification of character in a conversation}
The program loads separately each conversation and look for characters appearance into it. Character are always separated into speaker and listener. \\
\begin{itemize}
\item Firstly all dialogs of the conversation are loaded and the program extract characters from the potential speaker and listener and filter them using the procedure given in section \ref{filtering}. 

\item If no speaker is found, the program will look for sentences of the context located around the dialog and try to extract names from it. In some situations the speaker is mentioned just before or after the dialog, such that in the following sentence of Harry Potter 1: “Oh, yes,” said Mr. Dursley, his heart sinking horribly. “Yes, I quite agree.”. In this sentence the speaker, ``Mr Dursley'', is mentioned just between the line of dialogs. If no speaker are found, the program will also look in the context associated to the previous and next dialogs.\\

\item After this step, the number of appearance of each speaker in the conversation is counted. If a dialog have multiple speaker, only the speaker that have the most appearance in the conversation is kept.

\item In the last step, if the speaker of some dialogs stay unknown, we will try to infer the speaker from data of other dialogs.
	\begin{enumerate}
	\item If there is only one speaker in the conversation, we assume this is a monologue and all dialogs are assigned to this speaker.
	\item If there is 2 speakers, we assume that 2 characters are talking to each others. Firstly, the program will look for the speaker of the last dialog and assigned the other speaker to the dialog. If the conversation is beginning or the last dialog don't have an identified speaker, the program will look for the next dialog. Again, if a speaker is found, the other speaker is assigned to the dialog.
	\item If there is more than 2 speakers, the program will look for the speakers of the previous and next dialogs. Then it will assign to the dialog the most common speaker from the conversation among all detected speakers except the speakers previously identified.
	\end{enumerate}
\item Again, the number of appearance of each speaker in the conversation is counted. If a dialog have multiple speaker, only the speaker that have the most appearance in the conversation is kept.

\item Each sentence of context is linked with the speaker of the nearest dialog.

\item If a dialog have no identified listener, the program consider that all characters appearing in the conversation are listener, except the speaker of the dialog.
\end{itemize}
This method has the disadvantage of giving more appearance of most common speaker in a conversation. It also tend to add in conversation character that are mentioned without being part of it. For example if in Harry Potter, Harry told to Ron ``Don't say that to Hermionne'', she could wrongly be considered as the audience of the sentence. However in such case, the social network that will use this information would still be meaningful, as Harry speaking about Hermionne to Ron, implies a relationship between those 3 characters.


\subsubsection{Filtering of useless character}
All characters that don't have been identified a single time among all conversations are removed from the list of character. They will not be part of the social network.

\section{Construction of the network}
The social network is constructed using the identification of characters in dialogs in section \ref{identification}.
Two types of network are constructed:
\begin{enumerate}
\item context networks: social network that are relative to a single conversation and the associated context. 
\item an incremental network: a network which the result of the incremental addition of all context networks. This is a network representing all relations of characters in the fiction and information from all conversations. 
\end{enumerate}
\subsection{context networks}
Context network are a particular type of network developed by \cite{original}. They have the particularity of representing only a conversation, so, all characters of those networks tend to be linked between them. They have been builded following the intuition that the incremental addition of context networks can be used to learn about the evolution of a narration.\\

They are builded by executing the following steps:
\begin{enumerate}
\item A conversation and the related context are isolated.
\item For all dialogs having an identified speaker and audience, a link is added between the speaker and each of the members of the audience. The link is weighted following the number of dialog that links the couple of characters. This weight initialized at 1. There is also a second weight representing the sentiment between the characters. It is initialized with the value of the polarity detected in the line of dialog.
\item If there is already a link between a speaker and a listener, the first weight is incremented.  The polarity of the dialog is added to the sentiment-weight.
\item At the end of the network construction, the sentiment weight of each link is normalized by dividing it by the first weight of the link.
\end{enumerate}

\susection{incremental networks}
The name ``incremental network'' come from \cite{original} in opposition from the context networks. However incremental networks can be considered as standard social network extracted from fiction that have been studied in many article such as \cite{agarwal-etal-2013-automatic, NER, character_country, movie, fiction, Dekker2018EvaluatingSN}.
 


    %37/109


%\textit{Proper names detection}: For each chunk of each dialog context, if the main word is labelled as a proper noun, is composed of at least 3 letters and is not only uppercase, this noun is added in the alias table.\\
%\textit{Aliases detection}: All chunks are read again and the program compute the number of times that each pair of proper names appear together in a chunk. If 2 words appear more than 2/3 of the time together, they are considered as alias.




\bibliography{biblio} 
%\printbibliography%

\end{document}

% Template conçu par Benjamin Vanhemelryck et revu par François Bronchart - Mai 2013